\section{Remaining Work: Verified Compilation}
\label{sec:proposed-research}

% Source: Oliver's research plan and status update, November 20, 2026.
\sourcenote{New: proposed library and connection to the preceding projects.}
The preceding projects supply small proofs, effect-safe extraction,
  and support for proofs involving analyses and variable binding.
Using them in a verified compiler still requires a connection between
  the optimizer's certificates and the semantics of the programs it transforms.
The remaining work is a library that makes this connection reusable:
  developers supply a language and proved theorems in Lean or Rocq,
  and the library supplies an \texttt{egglog} optimizer whose results are checked at compiler runtime.

% Sources: Oliver's status update; ../extended-proofs/semantics/ENCODING.md,
% Proof status; EgglogSemantics/Proofs/SourceProof/Soundness.lean.
\sourcenote{New: completed checker and remaining integration, from Oliver's research plan.}
The Lean proof-checker prototype and its soundness proof for the supported fragment are complete.
The missing work is the round trip: exporting proved rules, running \texttt{egglog},
  and validating its returned program and certificate.
The initial implementation targets Lean, where the prototype exists;
  a Rocq application requires a corresponding verified interface.

\subsection{Background: Verified Compilation}
\label{sec:background-verification}

% Primary source: Leroy, A Formally Verified Compiler Back-end, JAR 2009.
\sourcenote{New background synthesis: compiler correctness and CompCert.}
A compiler correctness theorem states that successful compilation preserves
  the source program's observable behavior.
Proof assistants such as Lean and Rocq check this argument against formal language semantics.
CompCert is a verified compiler that composes correctness proofs for individual passes
  into a theorem relating source and target programs~\cite{compcert-backend}.

% Primary sources: Necula, PLDI 2000; Tristan and Leroy, PLDI 2009.
\sourcenote{New background synthesis: translation validation and verified validators.}
\emph{Translation validation} checks each compilation result against its input~\cite{translation-validation}.
The optimizer produces a candidate output, and a validator checks
  that the transformation preserves program behavior.
A verified validator has a soundness theorem establishing that every accepted transformation is correct.
Such validators already run as executable components of CompCert,
  allowing optimization heuristics to remain unverified~\cite{verified-lazy-code-motion}.

\subsection{A Library for Verified Optimization}
\label{sec:verified-compilation}
\label{sec:lean}

% Source: Oliver's November 20 research plan and clarification of the congruence
% relation and developer-supplied IR checks. This interface remains proposed.
\sourcenote{New proposed design: definitions and proofs supplied by the developer.}
The developer defines an intermediate representation and its semantics,
  and chooses a congruence relation $\simeq$ for the optimizer's equalities,
  with proofs of its laws.
This relation must be reflexive, symmetric, and transitive,
  and preserved by every term constructor.
These laws justify the equality reasoning performed by \texttt{egglog}:
  replacing a child by a related term produces a related enclosing term.
The developer registers theorems justifying rewrites and analysis rules,
  and supplies IR-specific validity checks with proofs of their soundness.

% Source: Oliver's proposed theorem-to-rule interface; formal shape proposed here.
\sourcenote{New proposed design: exporting registered theorems as rules.}
The library exports each supported rewrite theorem as an \texttt{egglog} rule,
  avoiding a handwritten rule maintained separately from its justification.
The initial interface supports conditional equalities of the form
\[
  \forall\vec{x},\quad
  P_1(\vec{x})\land\cdots\land P_n(\vec{x})
  \;\Longrightarrow\;
  \ell(\vec{x})\simeq r(\vec{x}),
\]
where $\vec{x}$ names the pattern variables and each $P_i$ is a precondition.
The exported rule matches $\ell$, requires evidence for the preconditions,
  and adds the equality with $r$.
For example, the summation theorem in Equation~\ref{eq:distribution-rule}
  exports a distribution rule guarded by $\mathtt{NotFree}$.
This interface supports a declared fragment of theorems;
  translating arbitrary Lean or Rocq propositions into rewrite rules is outside the initial scope.

% Source: Oliver's deployment plan; report Sections 3--5.
\sourcenote{New proposed design: compiler construction and runtime validation.}
The library also uses the registered theorems and IR-specific checks to instantiate an executable validator.
The proof assistant checks the theorems and the validator's correctness once,
  when the compiler is built (Figure~\ref{fig:verified-compilation-pipeline}, upper row).
For each input program, the compiler encodes the program for \texttt{egglog},
  runs search, and extracts a candidate (lower row).
Proof production supplies a certificate relating the input to that candidate.
The validator checks the certificate and IR-specific conditions,
  and the compiler accepts the candidate only if these checks succeed.
Here runtime means compiler execution, not execution of the generated program.
These per-input runs do not invoke the proof assistant.

% Source: Oliver's clarification of the verification boundary.
\sourcenote{New: benefit of using an unverified optimizer with a verified checker.}
\texttt{Egglog}'s implementation remains unverified.
It supplies candidate programs and certificates, which the verified checker validates
  against the registered theorems.
This allows a verified compiler to use \texttt{egglog}'s substantial body of unverified code,
  including its query engine, scheduling, and search optimizations,
  without proving that implementation correct.

% New diagram of Oliver's proposed framework; no source paper claims this design.
\begin{figure}[htbp]
  \sourcenote{New diagram of the proposed library and compiler workflow.}
  \centering
  % New diagram of Oliver's proposed verified-compilation library.
% Labels distinguish developer input, automatic work, and generated artifacts.
\begin{tikzpicture}[
    font=\footnotesize,
    stage/.style={draw, rounded corners=2pt, align=center, inner sep=4pt},
    build/.style={stage, text width=4cm, minimum height=2.45cm},
    run/.style={stage, text width=3.1cm, minimum height=2.05cm},
    flow/.style={-{Latex[length=2.2mm]}, thick}]
  \node[anchor=west, font=\small\bfseries] at (-1.65,1.65)
    {Once per compiler: in Lean or Rocq};
  \node[build, fill=orange!12] (theorems) at (0.45,0)
    {\textbf{Developer provides}\\[3pt]
     IR and semantics\\Congruence and laws\\
     Rewrite/analysis proofs\\IR validity check + proof};
  \node[build, fill=blue!8] (library) at (5.9,0)
    {\textbf{Library (automatic)}\\[3pt]
     Export rules\\Instantiate validator\\Derive pass theorem};
  \node[build, fill=black!4] (artifacts) at (11.35,0)
    {\textbf{Generated artifacts}\\[3pt]
     \texttt{egglog} rules\\Executable validator\\Correctness theorem};
  \draw[flow] (theorems) -- (library);
  \draw[flow] (library) -- (artifacts);
  \draw[black!30] (-1.65,-1.55) -- (13.45,-1.55);
  \node[anchor=west, font=\small\bfseries] at (-1.65,-1.95)
    {Each input program: automatic compilation without Lean/Rocq};
  \node[run, fill=orange!12] (input) at (0,-3.4)
    {\textbf{User provides}\\Input program};
  \node[run, fill=blue!8] (search) at (3.95,-3.4)
    {\textbf{Unverified}\\\texttt{egglog}\\Search + extraction\\Produce certificates};
  \node[run, fill=blue!8] (check) at (7.9,-3.4)
    {\textbf{Verified validator}\\Check certificate\\Check IR validity};
  \node[run, fill=black!4] (output) at (11.85,-3.4)
    {\textbf{Output}\\Accepted program};
  \draw[flow] (input) -- (search);
  \draw[flow] (search) -- (check);
  \draw[flow] (check) -- node[above, font=\footnotesize] {yes} (output);
  \draw[flow] (input.south) -- (0,-4.85)
    -- node[below, font=\footnotesize] {original input} (7.9,-4.85)
    -- (check.south);
\end{tikzpicture}

  \caption{Proposed division of work for a supported language interface.
    The compiler developer supplies the definitions and proofs at upper left;
    the library automatically exports rules, instantiates the validator,
    and derives the pass theorem.
    For each user-supplied input program, search, extraction, proof production,
    and validation are automatic.
    The generated rules guide the unverified \texttt{egglog} engine,
    and the validator combines certificate checking with checks supplied for the IR,
    such as effect safety.}
  \label{fig:verified-compilation-pipeline}
\end{figure}

\subsection{Correctness and the Round Trip}
\label{sec:verified-round-trip}

% Source: Oliver's requested compiler-runtime correctness theorem;
% notation and statement are proposed here.
\sourcenote{New proposed theorem: correctness of accepted compilation results.}
The library's main theorem connects acceptance by the validator to program semantics.
Fix a library $R$ of registered rules with their semantic proofs.
Let $p$ be a well-formed input, $q$ a candidate output, and $\pi$ its certificate.
Acceptance must establish that the two programs are related by the chosen congruence:
\begin{equation}
  \operatorname{validate}_{R}(p,q,\pi)=\mathsf{true}
  \quad\Longrightarrow\quad
  p\simeq q.
  \label{eq:validator-soundness}
\end{equation}
If checking fails, an optimization pass can return its unchanged input;
  a pass between different languages must instead report failure or use an independently justified fallback.
Connecting $\simeq$ to the language's semantics gives the compiler's behavior-preservation theorem.

% Source: Oliver's clarification that the semantic bridge needs investigation;
% existing proposal's requirement to support executions that stop before saturation.
\sourcenote{New open question: connecting the existing checker theorem to program semantics.}
The existing checker theorem concerns consequences of \texttt{egglog} rules
  in an e-graph closed under those rules (Section~\ref{sec:extension-correctness}).
Connecting that theorem to the semantic relation $\simeq$ requires further investigation.
The semantic argument must also hold when search stops at a time or size limit,
  without assuming that \texttt{egglog} reached saturation.
The proof and any changes needed to the formalization are part of the proposed work.

% Source: Oliver's clarification that analysis facts must be checked for specific
% terms, with the interpretation and implementation still open.
\sourcenote{New proposed work: interpreting and checking proofs of analysis facts.}
This semantic connection must also account for the analysis premises of conditional rewrites.
\texttt{Egglog}'s proof production from Section~\ref{sec:extension-proofs}
  supplies derivations of these facts as well as equalities.
The verified interface must establish each analysis premise for the specific terms
  used by the rewrite.
This part of the verified interface remains to be designed and implemented.

% Source: Oliver's proposed round trip; source-rule identity and export checks
% remain proposed features, not implemented parts of the prototype.
\sourcenote{New proposed design: correspondence between theorems and exported rules.}
The round trip must preserve the connection between exported rules and their registered theorems.
Each exported rule carries an identifier for its registered theorem.
A certificate names that rule, its substitution, and proofs of its premises.
The checker validates the rule instance and its premises against the registered rule.
The library must also justify the translation from theorem syntax into the checker's rule representation,
  for example by generating a correspondence proof checked when the optimizer is built.

% Source: Oliver's proposed integration; Section 5 checker scope.
\sourcenote{New proposed proof obligations: representation and certificate boundaries.}
The certified terms must also correspond to the compiler's input and output programs.
Encoding must preserve the input program's meaning;
  certificate import must preserve rule identities and substitutions;
  and decoding must connect the certified output term to the program actually returned.
The validator must reject malformed or ill-typed results and certificates for different endpoints.

% Source: Oliver's clarification that effect-safety checking is IR-specific.
\sourcenote{New proposed design: effect-safety checks supplied by the compiler developer.}
For an effectful IR, decoding also requires a valid ordering of effects.
The compiler developer supplies an effect-safety checker
  and a proof that it accepts only programs with valid effect ordering.
The runtime validator checks both the equality certificate and this IR-specific condition,
  allowing extraction itself to remain unverified.

% Source: proposed deployment boundary; no new trust-free execution claim.
\sourcenote{New proposed design: execution assumptions and deployment.}
The deployment argument must also account for how the verified validator becomes executable code.
This includes assumptions about code generation, the runtime, and communication with \texttt{egglog}.

\subsection{Candidate Case Studies}
\label{sec:verified-case-studies}

% Source: Oliver's three candidates; primary sources cited in each row.
% CompCert: https://compcert.org/doc/ (RTL, passes, and static analyses).
% NetKAT: Anderson et al., POPL 2014; KATch: Moeller et al., PLDI 2024.
% Luminal: https://github.com/luminal-ai/luminal and
% https://blog.luminal.com/p/compiling-models-to-megakernels.
% Effectful Luminal scope and NetKAT limitation clarified by Oliver after review.
\sourcenote{New comparison of Oliver's candidate case studies;
  existing systems from their papers and documentation, proposed pass scopes added.}
The case study must demonstrate both a correctness theorem and a useful optimization.
The candidates in Table~\ref{tab:compiler-cases} are alternatives for one main case study.
Selection depends on available semantics, integration effort, and useful optimizations
  that exercise analysis premises and, where applicable, effect-safe extraction.

\begin{table}[htbp]
  \centering
  \begin{tabular}{@{}p{0.17\linewidth}p{0.78\linewidth}@{}}
    \toprule
    Candidate & Proposed scope and main tradeoff \\
    \midrule
    CompCert &
    Optimize arithmetic and memory operations within a basic block
      (a straight-line instruction sequence) of its register-based RTL representation~\cite{compcert-doc}.
    This directly tests effectful optimization in a verified pipeline,
      but requires a Rocq checker interface and proofs connecting RTL to the extraction representation. \\[6pt]
    NetKAT &
    NetKAT is a language for packet-processing policies~\cite{netkat}.
    Simplify its policies by eliminating redundant tests and branches, initially excluding iteration.
    Its equational theory supports theorem export; KATch supplies benchmarks and an independent
      equivalence checker~\cite{katch}.
    Measure policy size and downstream compilation cost.
    Its drawback for this thesis is that it does not need effect-safe extraction. \\[6pt]
    Luminal &
    Luminal uses \texttt{egglog} to optimize tensor computations~\cite{luminal,luminal-megakernels}.
    Verify an effectful tensor-IR pass that needs effect-safe extraction,
      together with proofs of shape and bounds conditions.
    Compare overhead and generated-program performance on the same models.
    The theorem covers a selected pass, with explicit numerical semantics:
      real-number equalities do not justify arbitrary floating-point reassociation.
    GPU backend verification is separate. \\
    \bottomrule
  \end{tabular}
  \caption{Alternative case studies for the verified-optimization library.}
  \label{tab:compiler-cases}
\end{table}

\subsection{Evaluation Plan}
\label{sec:evaluation}

% Source: Oliver's goals for a reusable framework and a practical case study.
\sourcenote{New evaluation plan: correctness, automation, and developer effort.}
The evaluation asks whether the library makes useful optimizations easier to verify and deploy.
The formal deliverable is Equation~\ref{eq:validator-soundness}
  instantiated for the selected pass and composed with its compiler interface.
To assess reuse, the study adds a rewrite and an analysis rule after initial integration,
  measuring the additional definitions, proofs, and adapter code required.
These additions should reuse the generic checker proof and runtime integration.
Malformed certificates, missing premises, and mismatched endpoints provide negative integration tests.

% Source: Oliver's runtime-checking proposal and prior evaluation plan.
\sourcenote{New evaluation plan: runtime costs and optimization benefit.}
Measurements separate compiler construction from per-input costs:
  search, extraction, proof recording, reconstruction, minimization, and checking.
Runs with and without validation isolate its overhead;
  comparisons with proof-term reconstruction use the same transformations and semantic claims
  where both systems support them.
The study also reports certificate size, peak memory, validation failures,
  and the performance or code size of accepted outputs relative to the existing compiler.
Accepting unchanged programs establishes correctness but does not demonstrate optimization benefit.

% Source: existing report evaluation plan; Small Proofs and eggcc structural parameters.
\sourcenote{New evaluation plan for the parameterized-algorithm contribution.}
The evaluation records the structure that drives proof and extraction costs:
  proof-graph size, the number of congruence justifications, and statewalk width
  for workloads using effect-safe extraction.
Inputs of similar size but different parameter values test how well these measures explain runtime.
Varying the greedy proof-search budget measures the tradeoff between minimizing certificates
  and checking them.
Analyses and renaming introduce additional proof steps,
  so their counts are recorded separately rather than assuming that congruence alone explains the cost.
These experiments test whether the benefits from Sections~\ref{sec:proofs} and~\ref{sec:effectful-extraction}
  persist in the integrated compiler.

\subsection{Scope and Completion Criteria}
\label{sec:contexts}

% Source: Oliver's completion status and case-study alternatives;
% dependency ordering and scope are new proposal text.
\sourcenote{New research plan: remaining milestones and scope.}
The work proceeds through three milestones.
First, connect the existing Lean checker to \texttt{egglog} on a small language,
  completing theorem export and certificate import for a conditional rewrite whose premise comes from an analysis.
Second, package the interface as a library and establish the round-trip correctness theorem.
Third, integrate one case study and measure validation cost, optimization benefit, and extension effort.
The first milestone provides a concrete test of the interface before committing to a larger compiler.

% Source: proposed limits and contingency plan, based on Oliver's candidates.
\sourcenote{New research plan: limits and contingencies.}
The project will integrate the library with one proof assistant
  and evaluate it in a substantive compiler case study.
Supporting both Lean and Rocq, every \texttt{egglog} primitive,
  or extensions beyond those in Section~\ref{sec:extensions} is optional.
If a candidate's semantics makes integration too large, a smaller pass must still
  retain conditional rewrites and a useful optimization target.
If certificates dominate compilation time, minimization and sharing become evaluation priorities.
The evaluation must identify which parts of the thesis the chosen study exercises.

\subsection{Related Work}
\label{sec:related-verified-compilation}

% Primary sources: Tristan and Leroy, PLDI 2009, especially Sections 4 and 8;
% Rideau and Leroy, CC 2010.
\sourcenote{New literature synthesis: verified validation of lazy code motion
  and register allocation.}
Tristan and Leroy's CompCert validator checks lazy code motion,
  an optimization that moves computations to avoid repeated work~\cite{verified-lazy-code-motion}.
Its checks include whether moving a potentially failing computation changes program behavior.
Rideau and Leroy validate assignments of program values to registers or memory
  with a verified dataflow analysis~\cite{validated-register-allocation}.
These systems already separate optimization heuristics from correctness
  and check properties needed by effectful transformations.
Our proposed checker instead validates derivations assembled from registered rewrite and analysis rules,
  enabling reuse across optimizations.

% Primary sources: Tatlock and Lerner, PLDI 2010; Gross et al., ITP 2022.
\sourcenote{New literature synthesis: extensible verified compilers and theorem-driven rewriting.}
XCert extends CompCert with a verified engine for optimizations expressed in a dedicated language~\cite{xcert}.
Gross and colleagues compile proved Coq theorems into executable rewriters for Fiat Cryptography~\cite{verified-rewriting}.
Both reduce proof effort through reusable infrastructure.
Our proposed library applies that goal to \texttt{egglog}'s equality saturation and analyses,
  leaving search unverified and checking certificates at compiler runtime.
Its contribution depends on the supported rule interface, analysis premises,
  and practical integration and validation costs.

% Primary sources: Stepp, Tate, and Lerner, CAV 2011; Lopes et al., PLDI 2021.
\sourcenote{New literature synthesis: Peggy's LLVM validator and Alive2.}
Peggy's validator for LLVM's intermediate representation searches for equalities
  between original and optimized programs, including their memory operations~\cite{llvm-translation-validation}.
Alive2 uses satisfiability modulo theories (SMT) solvers to check that the transformed program
  introduces no behavior forbidden by the original, considering loops up to a bounded number of iterations~\cite{alive2}.
In our proposed compiler, the optimizer records a certificate during search,
  and a verified checker validates it after extraction.
The certificate must justify analysis-dependent rewrites,
  and the compiler's correctness argument must connect the certified equality to the executable output.

% Primary sources: Rossel et al., POPL 2026, Section 5;
% Bourgeat's dissertation, 2023, Section 7.2; Ekici et al., CAV 2017, Section 1.
\sourcenote{New literature synthesis: Lean's \texttt{egg} tactic,
  Rocq's Coquetier, and SMTCoq.}
Lean's \texttt{egg}~\cite{lean-egg} and Rocq's Coquetier~\cite{coquetier}
  are tactics, procedures that help construct proofs.
They reconstruct native proof terms from \texttt{egg}'s explanations.
SMTCoq instead evaluates a verified certificate checker inside Rocq
  and applies its soundness theorem~\cite{smtcoq}.
Our proposed library follows the verified-checker approach but runs the compiled validator
  during compilation, without constructing a new proof-assistant proof term for every input.
Performance comparisons must account for certificate production and checking
  as well as the tactics' proof reconstruction.

% Primary source: https://adam.chlipala.net/papers/EgraphsPOPL27/
% Only the authors' abstract is available; rechecked 2026-11-20.
\sourcenote{New literature synthesis: Jamner and colleagues' forthcoming work,
  based on the authors' published abstract.}
Jamner and colleagues' forthcoming work verifies an \texttt{egglog}-style
  equality-saturation procedure in Rocq, eliminating certificates~\cite{verified-egraphs}.
Our architecture leaves search unverified, allowing heuristics and schedules to change
  independently of the checker as long as they produce supported certificates.
Both approaches can yield executable tools; they place the proof obligation on different components.
